\documentclass[11pt]{article}
\newcommand{\doctitle}{Code Review: analysis/2025/pm\_functions.R}
\newcommand{\shorttitle}{Tidyverse Code Review}
\newcommand{\docdate}{2026-03-24}
% pmsimstats-ng standard document preamble
% Usage: % pmsimstats-ng standard document preamble
% Usage: % pmsimstats-ng standard document preamble
% Usage: \input{pmsimstats-preamble} after \documentclass[11pt]{article}
%
% Documents must define before \input:
%   \newcommand{\doctitle}{...}       Full document title
%   \newcommand{\shorttitle}{...}     Short title for header
%   \newcommand{\docdate}{YYYY-MM-DD} Document date

\usepackage[margin=1in]{geometry}
\usepackage{amsmath,amssymb}
\usepackage[T1]{fontenc}
\usepackage{lmodern}
\usepackage{microtype}
\usepackage{graphicx}
\graphicspath{{figures/}}
\usepackage{booktabs}
\usepackage{longtable}
\usepackage{array}
\usepackage{hyperref}
\usepackage{xcolor}
\usepackage{fancyhdr}
\usepackage{parskip}
\usepackage{caption}
\usepackage{enumitem}
\usepackage{verbatim}
\usepackage{listings}

\lstset{
  language=R,
  basicstyle=\small\ttfamily,
  keywordstyle=\color{blue!70!black},
  commentstyle=\color{green!50!black},
  stringstyle=\color{red!60!black},
  breaklines=true,
  frame=single,
  framerule=0.5pt,
  rulecolor=\color{gray!40},
  backgroundcolor=\color{gray!5},
  xleftmargin=1em,
  framexleftmargin=1em,
  columns=flexible
}

\hypersetup{
  colorlinks=true,
  linkcolor=blue!60!black,
  citecolor=green!50!black,
  urlcolor=blue!60!black
}

\pagestyle{fancy}
\fancyhf{}
\lhead{\footnotesize pmsimstats-ng Technical Report}
\rhead{\footnotesize \shorttitle}
\rfoot{\footnotesize \thepage}
\renewcommand{\headrulewidth}{0.4pt}

\title{\doctitle}
\author{pmsimstats team}
\date{\docdate}
 after \documentclass[11pt]{article}
%
% Documents must define before \input:
%   \newcommand{\doctitle}{...}       Full document title
%   \newcommand{\shorttitle}{...}     Short title for header
%   \newcommand{\docdate}{YYYY-MM-DD} Document date

\usepackage[margin=1in]{geometry}
\usepackage{amsmath,amssymb}
\usepackage[T1]{fontenc}
\usepackage{lmodern}
\usepackage{microtype}
\usepackage{graphicx}
\graphicspath{{figures/}}
\usepackage{booktabs}
\usepackage{longtable}
\usepackage{array}
\usepackage{hyperref}
\usepackage{xcolor}
\usepackage{fancyhdr}
\usepackage{parskip}
\usepackage{caption}
\usepackage{enumitem}
\usepackage{verbatim}
\usepackage{listings}

\lstset{
  language=R,
  basicstyle=\small\ttfamily,
  keywordstyle=\color{blue!70!black},
  commentstyle=\color{green!50!black},
  stringstyle=\color{red!60!black},
  breaklines=true,
  frame=single,
  framerule=0.5pt,
  rulecolor=\color{gray!40},
  backgroundcolor=\color{gray!5},
  xleftmargin=1em,
  framexleftmargin=1em,
  columns=flexible
}

\hypersetup{
  colorlinks=true,
  linkcolor=blue!60!black,
  citecolor=green!50!black,
  urlcolor=blue!60!black
}

\pagestyle{fancy}
\fancyhf{}
\lhead{\footnotesize pmsimstats-ng Technical Report}
\rhead{\footnotesize \shorttitle}
\rfoot{\footnotesize \thepage}
\renewcommand{\headrulewidth}{0.4pt}

\title{\doctitle}
\author{pmsimstats team}
\date{\docdate}
 after \documentclass[11pt]{article}
%
% Documents must define before \input:
%   \newcommand{\doctitle}{...}       Full document title
%   \newcommand{\shorttitle}{...}     Short title for header
%   \newcommand{\docdate}{YYYY-MM-DD} Document date

\usepackage[margin=1in]{geometry}
\usepackage{amsmath,amssymb}
\usepackage[T1]{fontenc}
\usepackage{lmodern}
\usepackage{microtype}
\usepackage{graphicx}
\graphicspath{{figures/}}
\usepackage{booktabs}
\usepackage{longtable}
\usepackage{array}
\usepackage{hyperref}
\usepackage{xcolor}
\usepackage{fancyhdr}
\usepackage{parskip}
\usepackage{caption}
\usepackage{enumitem}
\usepackage{verbatim}
\usepackage{listings}

\lstset{
  language=R,
  basicstyle=\small\ttfamily,
  keywordstyle=\color{blue!70!black},
  commentstyle=\color{green!50!black},
  stringstyle=\color{red!60!black},
  breaklines=true,
  frame=single,
  framerule=0.5pt,
  rulecolor=\color{gray!40},
  backgroundcolor=\color{gray!5},
  xleftmargin=1em,
  framexleftmargin=1em,
  columns=flexible
}

\hypersetup{
  colorlinks=true,
  linkcolor=blue!60!black,
  citecolor=green!50!black,
  urlcolor=blue!60!black
}

\pagestyle{fancy}
\fancyhf{}
\lhead{\footnotesize pmsimstats-ng Technical Report}
\rhead{\footnotesize \shorttitle}
\rfoot{\footnotesize \thepage}
\renewcommand{\headrulewidth}{0.4pt}

\title{\doctitle}
\author{pmsimstats team}
\date{\docdate}


\begin{document}
\maketitle
\thispagestyle{fancy}
\tableofcontents

\section{Dead Code}

\subsection{1. \texttt{calculate\_carryover()} (lines 24-77)}

Multi-model carryover function supporting exponential, linear, and
Weibull decay. \textbf{Never called} by any function in the file. The
actual carryover is computed inline in \texttt{apply\_carryover\_to\_component()}
using the exponential formula \texttt{(1/2)\^{}(tsd/halflife)}. The linear
and Weibull models were never integrated into the DGP.

\textbf{Action:} Remove entirely.

\subsection{2. \texttt{calculate\_component\_halflives()} (lines 138-145)}

Returns a list with three identical values (\texttt{br\_halflife},
\texttt{pb\_halflife}, \texttt{tv\_halflife}), all set to \texttt{base\_halflife}. Only
\texttt{br\_halflife} is ever used (the function guards on
\texttt{component\_name != "br"} in \texttt{apply\_carryover\_to\_component()}).

\textbf{Action:} Inline the single value. Replace all calls with
\texttt{model\_param\$carryover\_t1half} directly.

\subsection{3. \texttt{calculate\_bio\_response\_with\_interaction()} (lines 95-126)}

Thin wrapper that:

\begin{enumerate}
\item Calls \texttt{mod\_gompertz()} on \texttt{trial\_data\$tod}
\item Sets \texttt{bio\_response\_test <- base\_bio\_response == 0}
\item Calls \texttt{apply\_carryover\_to\_component()}
\item Returns a 3-element list
\end{enumerate}

This adds indirection without value. The orig code does this inline
in \texttt{buildSigma()} in 10 lines. The \texttt{raw\_bio\_response\_means} return
field is never used by any caller.

\textbf{Action:} Inline into \texttt{generate\_data()} and \texttt{build\_sigma\_matrix()}.

\subsection{4. \texttt{validate\_correlation\_structure()} (lines 778-813)}

Diagnostic function. Calls \texttt{build\_sigma\_matrix()} and reports
eigenvalues and condition number. Not called by the simulation
pipeline. Was useful during development but is now redundant with
\texttt{validate\_parameter\_grid()}.

\textbf{Action:} Remove. The orig's \texttt{validateParameterGrid()} (ported
as Step 7c) covers this use case.

\subsection{5. \texttt{create\_sigma\_cache\_key()} (lines 815-822)}

Creates a cache key from \texttt{design\_name}, \texttt{params\$n\_participants},
\texttt{params\$biomarker\_correlation}, \texttt{params\$carryover\_t1half}. \textbf{Never
called.} The actual caching in \texttt{generate\_simulated\_results()} uses
\texttt{paste(vpg indices, sep='\_')} directly.

\textbf{Action:} Remove entirely.

\subsection{6. \texttt{validate\_parameter\_grid()} (lines 828-991)}

170-line validation function with extensive console output, emoji,
condition number statistics, and recommendations. References
\texttt{param\_grid\$biomarker\_correlation} and \texttt{param\_grid\$n\_participants}
-- column names from the old 2025 parameter format that no longer
exists after the alignment. Would crash if called.

\textbf{Action:} Remove. The orig's \texttt{validateParameterGrid()} and the
ported \texttt{generate\_simulated\_results()} handle pre-flight validation.

\subsection{7. \texttt{report\_parameter\_validation()} (lines 997-1044)}

Companion reporting function for \texttt{validate\_parameter\_grid()}. Same
issue -- references old column names. Never called by the pipeline.

\textbf{Action:} Remove.

\subsection{8. Deprecated docstring block (lines 147-158)}

Roxygen comment block for a function that was already removed.
Describes the deprecated \texttt{calculate\_carryover\_adjusted\_correlations()}
which no longer exists.

\textbf{Action:} Remove.

\section{Duplication}

\subsection{9. \texttt{generate\_data()} and \texttt{build\_sigma\_matrix()} share \textasciitilde{}100 lines}

Both functions independently compute: labels, standard\_deviations,
means (including the Gompertz + carryover loop), and the correlation
matrix. The logic is copy-pasted. In orig, \texttt{buildSigma()} is a
standalone function and \texttt{generateData()} calls it (or uses a
cached version).

\textbf{Action:} Make \texttt{generate\_data()} always call \texttt{build\_sigma\_matrix()}
when no cached sigma is provided, eliminating the duplicated sigma
construction path (lines 433-582 of generate\_data). This matches
the orig architecture where \texttt{generateData()} calls \texttt{buildSigma()}.

\subsection{10. \texttt{factor\_types} and \texttt{factor\_abbreviations} are identical}

After the alignment, \texttt{factor\_types = c("tv", "pb", "br")} and
\texttt{factor\_abbreviations = c("tv", "pb", "br")} are the same vector.
Both are passed around as separate arguments and iterated in
parallel. The \texttt{current\_factor} / \texttt{factor\_abbrev} distinction no
longer exists.

\textbf{Action:} Collapse to a single \texttt{factors <- c("tv", "pb", "br")}
throughout.

\section{Simplifications}

\subsection{11. \texttt{apply\_carryover\_to\_component()} over-engineered}

Takes \texttt{component\_halflives} (a list) and \texttt{component\_name} (a
string), then immediately returns unchanged if \texttt{component\_name !=
"br"}. This function is called three times in a loop -- twice it
returns immediately. The function could simply accept the halflife
value directly and only be called for BR.

\textbf{Action:} Call only for BR component. Pass
\texttt{model\_param\$carryover\_t1half} directly instead of the halflives
list. Remove the \texttt{component\_name} guard.

\subsection{12. \texttt{prepare\_trial\_data()} fallback logic}

Lines 238-246 handle a \texttt{week} column fallback that is never
exercised. All trial designs in this codebase use \texttt{t\_wk}. The
\texttt{on\_drug} column it adds duplicates information already available
from \texttt{tod > 0}.

\textbf{Action:} Simplify to: add \texttt{t\_wk\_cumulative <- cumsum(t\_wk)} and
\texttt{on\_drug <- (tod > 0)}. Remove the \texttt{week} fallback.

\subsection{13. \texttt{build\_correlation\_matrix()} signature has 11 parameters}

Several are vestigial: \texttt{bio\_response\_test}, \texttt{bio\_response\_means},
\texttt{means} are only used for the BM-BR correlation block, and even
there \texttt{bio\_response\_test} is only checked when \texttt{trial\_data} is
NULL (which never happens in practice since \texttt{trial\_data} is always
passed). The \texttt{factor\_types} parameter is now identical to
\texttt{factor\_abbreviations}.

\textbf{Action:} Simplify signature. Remove \texttt{bio\_response\_test},
\texttt{bio\_response\_means}, \texttt{means}, \texttt{factor\_types}. Use \texttt{trial\_data}
(always present) for the on-drug check.

\subsection{14. \texttt{process\_participant\_data()} pre-allocates then overwrites}

Lines 1054-1058 create a zero-filled tibble and bind it, then
lines 1060-1079 overwrite every column. The pre-allocation has no
benefit since \texttt{mutate()} creates new columns regardless.

\textbf{Action:} Remove the pre-allocation block. Let the \texttt{mutate()}
calls create the columns directly.

\subsection{15. PD tracking in \texttt{generate\_data()} is single-call accounting}

\texttt{sigma\_count} and \texttt{non\_positive\_definite\_count} are initialized to
0, incremented at most once (the function builds one sigma), and
then attached as attributes. The \texttt{non\_positive\_definite\_rate} is
always 0 or 1. This tracking made sense in a loop but is trivial
for a single-sigma function.

\textbf{Action:} Replace with a single boolean attribute
\texttt{was\_pd\_corrected}.

\section{Estimated Impact}

\begin{table}[htbp]
\centering
\begin{tabular}{lrccc}
\toprule
Category & Items & Lines removed & Lines simplified \\
\midrule
Dead code & 7 & \textasciitilde{}350 & -- \\
Duplication & 2 & \textasciitilde{}100 & \textasciitilde{}50 \\
Simplifications & 5 & \textasciitilde{}30 & \textasciitilde{}60 \\
\midrule
\textbf{Total} & \textbf{14} & \textbf{\textasciitilde{}480} & \textbf{\textasciitilde{}110} \\
\bottomrule
\end{tabular}
\end{table}

Current file: 1650 lines. After cleanup: \textasciitilde{}1060 lines.

\vfill
\noindent\rule{\textwidth}{0.4pt}
{\footnotesize
Rendered on 2026-03-25 at 19:01 PDT.\\
Source: \textasciitilde/prj/alz/10-pmsimstats-ng/pmsimstats-ng/docs/16-tidyverse-code-review.tex}
\end{document}
